\documentclass{article}

% ====== Packages ==========
\usepackage[utf8]{inputenc}       % encoding
\usepackage[T1]{fontenc}          % better fonts
\usepackage{lmodern}              % Latin Modern font
\usepackage{graphicx}             % for including images
\usepackage{amsmath, amssymb}     % math symbols
\usepackage{hyperref}             % clickable links
\usepackage{geometry}            % adjust page margins
\usepackage{cite}                % improved citations

\geometry{margin=1in}

% ======== Title ==========
\title{Hybrid Automata-Based Software Architecture and Simulation Tools for USV Mission Planning}
\author{
    Ryan McKee$^{1}$ \and Nikolaos Athanasopoulos$^{2}$ \and Wasif Naeem$^{3}$ \\
    \small $^{1,3}$Queen's University Belfast \\
    \small $^{2}$Artemis Technologies \\
    \small \texttt{\{r.mckee,  n.athanasopoulos, w.naeem\}@qub.ac.uk}
}
\date{\today}

\begin{document}

\maketitle

\begin{abstract}
Modern maritime operations require autonomous navigation solutions that can adapt to dynamic environments and evolving regulations while maintaining compliance with COLREGs~\cite{colregs1972} and offering interpretable, certifiable safety. Deep Reinforcement Learning (DRL) has recently shown promise for autonomous collision avoidance (e.g., \cite{wang2022colregs}), but DRL agents are often opaque and can behave unpredictably outside their training distribution, complicating certification and operator trust.

This paper is part of a broader research programme on autonomous surface vessel mission planning and control, comprising several complementary studies on hybrid decision-making, risk assessment, and system integration. Within this context, we present a hybrid-automaton framework that combines the decision-making capabilities of advanced learning-based controllers with the transparency and verifiability of formal state-based control. Our formalism structures behavior into interpretable modes, enables rule-driven switching consistent with COLREGs, and supports the integration of DRL policies as mode controllers.

To enable reproducible development and deployment, we introduce \textbf{HybrautNav}, a modular ROS~2-based navigation stack whose tactical layer is defined by a hybrid automaton using the AMDL specification language. HybrautNav is packaged as a set of Docker-deployable modules suitable for embedded platforms. Beyond the navigation stack, we also present an overall software architecture that integrates HybrautNav into a broader mission-planning system with real-time risk assessment and scenario management. We evaluate the system in a configurable USV simulation suite and on representative encounter scenarios, measuring collision rates, COLREGs-compliance metrics, and operator interpretability. Results demonstrate that our hybrid approach retains the performance benefits of DRL while improving transparency and enabling easier regulatory assessment. The architecture and example scenarios are released to support community adoption and future research.
\end{abstract}

\section{Introduction}
Autonomous Surface Vessels (ASVs) also referred to as Marine Autonomous Surface Ships or simply MASS have shown great promise and potential to be disruptive in terms of technological, environmental and safety challenges at sea. With over $90\%$ of world trade carried out by sea, safety of vessels and onboard crew as well as their environmental footprint is not compromised. ASVs must operate safely and efficiently in dynamic maritime environments, where compliance with the International Regulations for Preventing Collisions at Sea (COLREGs)~\cite{colregs1972} is essential. While learning-based approaches such as Deep Reinforcement Learning (DRL) have shown promise for collision avoidance, their black-box nature complicates certification and can lead to unpredictable behavior in unseen scenarios~\cite{wang2022colregs}.

To address these challenges, this work proposes a hybrid automata-based navigation architecture that combines interpretable decision logic with advanced control strategies. Specifically, we introduce \textbf{HybrautNav}, a modular ROS~2 navigation stack whose tactical layer is defined by hybrid automata specified using the Automaton Model Definition Language (AMDL). This architecture is complemented by a mission-level risk assessment system that provides situational awareness and regulatory context for COLAV decisions.

\section{Related Work}
Eriksen et al.~\cite{eriksen2020hybrid} proposed a three-layer hybrid collision-avoidance (COLAV) system, separating high-level planning, mid-level COLREGs-aware MPC, and short-term reactive control. Their approach demonstrated successful collision avoidance in all tested scenarios, highlighting the effectiveness of a modular hybrid design for separating planning horizons (strategic $\rightarrow$ tactical $\rightarrow$ immediate). However, their work remained at the algorithmic and simulation level: no software architecture, middleware integration, or standardized deployment framework was provided.

- lately there is a focus on restricted waters collision avoidance
- waters with mixed traffic i.e. crewed and uncrewed vessels
- optimization based methods are some of the most utilised strategies
- COLREGs is not considered in many papers
- Hybrid collision avoidance is fairly common.

Building on this foundation, our work focuses on creating a reusable and extensible \textbf{software architecture} that implements similar hybrid principles within ROS~2. To our knowledge, no open-source navigation stack exists that integrates hybrid automata-based tactical decision-making with modular strategic and control layers for COLREGs-compliant USV navigation. Our contribution addresses this gap by introducing HybrautNav and demonstrating its integration within a complete mission-planning and risk assessment framework.

\section{Methodology}
The proposed system consists of several layers and modules designed to interoperate in real time. Figure~\ref{fig:architecture} shows the overall structure.

\subsection{System Overview}
\begin{itemize}
    \item \textbf{Strategic Layer}: Mission-level planning, risk assessment, and waypoint management.
    \item \textbf{Tactical Layer}: Hybrid automaton decision logic defined in AMDL, integrating DRL or classical controllers as mode-specific behaviors.
    \item \textbf{Control Layer}: Low-level COLAV controllers (e.g., MPC or PID) executing commands from the tactical layer.
\end{itemize}

%% TODO: Add High Level Diagrams

\subsection{HybrautNav Navigation Stack}
HybrautNav is implemented as a set of ROS~2 nodes and Dockerized modules. Each layer communicates using standard ROS interfaces, enabling easy swapping and testing of different planning or control algorithms. The tactical layer is implemented as a hybrid automaton node, with mode transitions triggered by COLREGs encounter classification and dynamic risk assessments.

\subsection{Hybrid Automaton Definition (AMDL)}
We specify tactical decision logic using the Automaton Modeling Definition Language (AMDL). Each automaton defines:
\begin{itemize}
    \item \textbf{Modes}: e.g., normal transit, overtaking, crossing, head-on avoidance, emergency.
    \item \textbf{Guards}: CPA/TCPA thresholds, bearing sectors, or regulatory conditions.
    \item \textbf{Actions}: Controller selection, waypoint updates, or COLAV maneuvers.
\end{itemize}

\subsection{Risk Assessment Module}
A mission-level risk assessment module continuously evaluates vessel trajectories, environmental data, and regulatory zones to provide a situational context for tactical decisions. This module can trigger mode transitions or alter strategic plans in response to evolving scenarios.

%% 

\section{Running Code}
%% TODO: Show how to run code, maybe examples of some diagrams


\section{Results}
We evaluated the system in a configurable simulation environment with multiple dynamic obstacles, both cooperative and non-compliant. Metrics included collision rate, COLREGs compliance score, and system interpretability (e.g., mode traceability and state transitions). The hybrid automaton tactical layer successfully enforced rule-based behaviors while allowing for learning-based control within modes.

\section{High Level Architecture}

%% Show how to link to different components, like gazebo simulation, foxglove, rviz, and so on. 

\section{Discussion}
The results indicate that a hybrid automaton tactical layer offers a strong balance between performance and interpretability. By structuring decision logic formally, the system is more amenable to verification and certification processes than pure DRL-based approaches. The ROS~2-based modular design also supports rapid prototyping and deployment across different USV platforms.

\section{Conclusion and Future Work}
This paper introduced HybrautNav, a hybrid automata-based navigation stack for USVs, integrated within a broader mission-planning and risk assessment system. Future work will focus on field trials, integration with sensor fusion pipelines, and extending AMDL specifications to cover more complex multi-agent encounters.

\bibliographystyle{plain}
\bibliography{references}

\end{document}
